% Generated draft: requires mathematical and source review.
\documentclass[11pt]{article}
\usepackage[margin=25mm]{geometry}
\usepackage{fontspec}
\setmainfont{Latin Modern Roman}
\usepackage{amsmath,amssymb,mathtools,unicode-math}
\setmathfont{Latin Modern Math}
\usepackage{enumitem,needspace,xurl,hyperref}
\setlength{\parindent}{0pt}
\setlength{\parskip}{5pt}
\setlength{\emergencystretch}{3em}
\newcommand{\N}{\mathbb{N}}
\newcommand{\Z}{\mathbb{Z}}
\newcommand{\Q}{\mathbb{Q}}
\newcommand{\R}{\mathbb{R}}
\newcommand{\C}{\mathbb{C}}
\newcommand{\F}{\mathbb{F}}
\newcommand{\A}{\mathbb{A}}
\newcommand{\PP}{\mathbb{P}}
\newcommand{\E}{\mathbb{E}}
\newcommand{\eps}{\varepsilon}
\newcommand{\dd}{\,\mathrm{d}}
\begin{document}
\section*{UnsolvedMath ID: 30006461}
\textbf{Source problem code:} catalogue entry
\subsection{Definitions and mathematical statement}
Let \(f_{\mathrm{Mult}}\) denote the limiting profile referred to in the supplied record for the multiplication-table problem. Its definition, domain, normalization, and governing asymptotic formula are not supplied. Write \(D\) for its intended domain, without specifying that domain. The supplied nonconstancy conjecture is
\[
\exists x,y\in D\quad f_{\mathrm{Mult}}(x)\ne f_{\mathrm{Mult}}(y).
\]
The record expresses the conjecture as
\[
\frac{\max_{x\in D} f_{\mathrm{Mult}}(x)}{\min_{x\in D} f_{\mathrm{Mult}}(x)}>1.
\]
This ratio formulation is equivalent to nonconstancy if the maximum and minimum exist as finite real numbers and the minimum is strictly positive. These hypotheses are not established in the supplied excerpt. Thus the display records the intended conjecture, but a complete mathematical formulation requires the missing definition of the profile and the hypotheses justifying its extrema and positivity.

For context, for a positive integer \(n\), put \([n]=\{1,\ldots,n\}\). The description of the multiplication table gives the intended counting function
\[
M(n)=\bigl|\{ab:a,b\in[n]\}\bigr|.
\]
The extracted source instead restricts the counted products to \([n]\), while calling the result the number of distinct entries in the \(n\times n\) table. These are incompatible descriptions; the displayed definition follows the verbal description and explicitly records this discrepancy. The source also describes \(M(n)\) as roughly the number of integers in \([n^2]\) having a divisor in \([n/2,n]\); it supplies no precise meaning of ``roughly.''

With logarithms interpreted as natural logarithms, the source gives the Erd\H{o}s--Tenenbaum--Ford constant
\[
\delta=1-\frac{1+\log\log 2}{\log 2}\approx0.0860713.
\]
It attributes the exponent-scale estimate \(M(n)=n^2(\log n)^{-\delta+o(1)}\) to Erd\H{o}s and the sharper comparison
\[
M(n)\asymp n^2(\log n)^{-\delta}(\log\log n)^{-3/2}
\qquad(n\to\infty)
\]
to Ford. Here \(A(n)\asymp B(n)\) means that positive constants \(c,C\) satisfy \(cB(n)\le A(n)\le CB(n)\) for all sufficiently large \(n\). These are supplied background attributions, not independently checked results.

The extracted passage appears to report
\[
\frac{\max_x f_{\mathrm{Mult}}(x)}{\min_x f_{\mathrm{Mult}}(x)}
\le \frac{\max_x g(x)}{\min_x g(x)}
<1+2\cdot10^{-7},
\]
but its formula layout is damaged and the definition of \(g\) is unreadable. This comparison cannot be certified from the supplied text. The supplied dated assessment attributes an asymptotic and a rigorous upper ratio bound to Green and Sawhney, while describing the strict lower bound above as unproved. The excerpt also says that an inequality of this type can be proved for divisors of smooth numbers in sufficiently short dyadic intervals; it supplies no quantitative hypotheses or precise theorem for that comparison. Its statement of ``Theorem 1 (Green-S. '25+)'' is truncated. None of these incomplete passages supplies the missing definition of \(f_{\mathrm{Mult}}\).
\subsection{Short English statement}
The intended conjecture says that the limiting profile for the number of distinct products in a multiplication table takes at least two different values. The supplied record does not define that profile or its domain. Its maximum-to-minimum ratio formulation additionally requires attained finite extrema and a strictly positive minimum. This is a statement-only draft requiring human mathematical review.
\subsection{Sources}
\begin{itemize}
\item Local UnsolvedMath registry, permanent ID \(30006461\), title ``Nonconstancy of the Multiplication-Table Limit Profile,'' exact source problem code \texttt{catalogue entry}. The supplied record lists proposed year 2025 and names no proposer. Its labels \texttt{open} and \texttt{corrected\_verified} are supplied registry classifications, not independent verification.
\item Supplied report citation: \emph{Analytic Number Theory (2026)}, Oberwolfach Reports, DOI 10.4171/owr/2025/51, pp.\ 2701--2764. \url{https://doi.org/10.4171/owr/2025/51}. The literature-source label is ``Oberwolfach report 51/2025.'' The supplied title year and report identifier use different years; their relationship has not been checked.
\item Supplied literature source: ``Green--Sawhney, permutation model.'' \url{https://arxiv.org/abs/2604.28116}. The record attributes the asymptotic and upper ratio bound to Green and Sawhney. No fuller bibliographic metadata is supplied here.
\item Literature assessment checked at \texttt{2026-08-22}: the supplied curation labels the problem open and says that Green and Sawhney obtain the asymptotic and a rigorous upper ratio bound but leave the strict nonconstancy inequality unproved. This is supplied dated curation derived from the OWR record, not an independently verified literature assessment. The record itself requests independent verification of time-sensitive status and cited claims.
\item The extracted background attributes earlier multiplication-table estimates to Erd\H{o}s, with source-local reference label [1], and Ford, with source-local reference label [5]. Their bibliography entries are absent from the supplied record, so these references cannot be identified more precisely here.
\end{itemize}
\end{document}
